\documentclass[10pt]{article}
\usepackage[margin=1in]{geometry}
\usepackage[T1]{fontenc}
\usepackage{lmodern,amsmath,amssymb,amsthm,microtype}
\usepackage[colorlinks=true,allcolors=blue]{hyperref}
\newtheorem{theorem}{Theorem}
\newtheorem{lemma}{Lemma}
\newtheorem{proposition}{Proposition}
\newcommand{\Z}{\mathbb Z}
\newcommand{\N}{\mathbb N}
\title{An alternating antimorphic Fine--Wilf theorem}
\author{Alec Kriebel\\\small Independent researcher\\
  \small\href{https://orcid.org/0009-0001-9320-500X}{ORCID: 0009-0001-9320-500X}}
\date{3 October 2026}
\begin{document}
\maketitle
\begin{abstract}
Let $\theta$ be an antimorphic involution of a free monoid. If a finite
word $w$ has alternating $\theta$-periods $p,q$ and
$|w|\ge p+q-\gcd(p,q)$, then $\gcd(p,q)$ is an alternating
$\theta$-period of $w$. This proves the unnumbered conjecture stated
after Theorem 23 in Nowotka's contribution, joint with Bischoff, to
\emph{Mini-Workshop: Combinatorics on Words} (2010). The proof applies
the classical Fine--Wilf theorem to the common reflected extension
$\theta(w)w$ and recovers the alternating phase at its center.
The length formula cannot be decreased by one uniformly.
\end{abstract}

\section{The statement}
Let $A$ be an arbitrary alphabet, and let $A^*$ be its free monoid,
with empty word $\varepsilon$. An \emph{antimorphic involution} is a map
$\theta:A^*\to A^*$ such that $\theta(xy)=\theta(y)\theta(x)$ and
$\theta^2=\mathrm{id}$. A positive integer $p$ is an
\emph{alternating $\theta$-period} of a finite word $w$ if there is a
word $u\in A^p$ such that $w$ is a prefix of $(u\theta(u))^\omega$.
The periods need not be minimal. This definition prescribes alternating
blocks; it differs from a prefix of an arbitrary concatenation of $u$
and $\theta(u)$.

\begin{theorem}\label{thm:main}
Suppose that $w$ has alternating $\theta$-periods $p$ and $q$, and put
$d=\gcd(p,q)$. If $L=|w|\ge p+q-d$, then, for the length-$d$ prefix
$v$ of $w$, the word $w$ is a prefix of $(v\theta(v))^\omega$.
\end{theorem}

This is the exact alternating conclusion in the unnumbered conjecture
following Theorem 23 on printed page 2220 of \cite{OWR}; it is distinct
from numbered Conjecture 27 there. Theorem 23 gives the sufficient
bound $p+q$. Bischoff's thesis \cite[Lemma 3.9 and Satz 3.11]{Bischoff}
also develops reflection and doubled ordinary periods and proves that
longer bound. We use those classical symmetries with a common extension
of the entire finite prefix. Fixed letters of $\theta$ are allowed;
in particular, ordinary reversal is included.

\section{The common extension}
For a word of length $n$, an \emph{ordinary period} $h>0$ means that
letters in positions $i$ and $i+h$ agree whenever both positions occur.
We use the classical Fine--Wilf theorem \cite{FW}: a word with ordinary
periods $h,k$ and length at least $h+k-\gcd(h,k)$ has ordinary period
$\gcd(h,k)$.

\begin{lemma}\label{lem:letters}
There is an involutive permutation $\sigma:A\to A$ such that
\[
\theta(a_0\cdots a_{n-1})=\sigma(a_{n-1})\cdots\sigma(a_0).
\]
\end{lemma}
\begin{proof}
The identity $\theta(\varepsilon)=\theta(\varepsilon)^2$ forces
$\theta(\varepsilon)=\varepsilon$. Since $\theta$ is bijective,
nonempty words have nonempty images. If the image of a letter were
$xy$ with both factors nonempty, applying $\theta$ would factor that
letter as the two nonempty words $\theta(y)\theta(x)$, a contradiction.
Thus letters map to letters. Involutivity and the antimorphism rule
give the assertion.
\end{proof}

\begin{lemma}\label{lem:extension}
If $L=|w|\ge p$ and $p$ is an alternating $\theta$-period of $w$,
then $\theta(w)w$ has ordinary period $2p$.
\end{lemma}
\begin{proof}
Write $u=w_0\cdots w_{p-1}$ and extend the block $u\theta(u)$
periodically to a sequence $(s_i)_{i\in\Z}$ of period $2p$, with
$s_i=w_i$ for $0\le i<L$. If $r$ is the residue of $i$ in
$\{0,\ldots,2p-1\}$, then
\[
 s_i=\begin{cases}u_r,&0\le r<p,\\
       \sigma(u_{2p-1-r}),&p\le r<2p.
       \end{cases}
\]
The residue of $-1-i$ is $2p-1-r$, so
$s_{-1-i}=\sigma(s_i)$ in both cases. Hence, on the negative positions
$-L\le j<0$, we have $s_j=\sigma(w_{-1-j})$.
The restriction to $[-L,L-1]$ is therefore exactly $\theta(w)w$,
indexed with $w_0$ at position $0$. As a restriction of a $2p$-periodic
sequence, it has ordinary period $2p$.
\end{proof}

The extension in Lemma~\ref{lem:extension} is determined by $w$ alone.
In particular, applying it with $p$ and $q$ yields ordinary periods
on the \emph{same finite word}. No equality of two unseen infinite
extensions is being assumed, and no condition on the last incomplete
block of $w$ is needed.

\begin{proof}[Proof of Theorem~\ref{thm:main}]
Since $d\le\min(p,q)$, the hypothesis gives $L\ge\max(p,q)$.
Lemma~\ref{lem:extension} shows that the word $W=\theta(w)w$, of
length $2L$, has ordinary periods $2p$ and $2q$. Moreover,
\[
 2L\ge 2p+2q-2d,\qquad \gcd(2p,2q)=2d.
\]
Fine--Wilf therefore gives ordinary period $2d$ for $W$.
Continue to index $W$ on $[-L,L-1]$. Its central restriction to
$[-d,d-1]$ is $\theta(v)v$, where $v=w_0\cdots w_{d-1}$.

For $0\le i<L$, let $r\in\{0,\ldots,2d-1\}$ be the residue
of $i$ modulo $2d$. If $r<d$, choose $j=r$. Otherwise choose
$j=r-2d\in[-d,-1]$. In each case $j$ and $i$ are in $[-L,L-1]$
and differ by a multiple of $2d$. All intermediate positions between
them remain in that interval, so ordinary period $2d$ gives $W_i=W_j$.
Consequently
\[
 w_i=\begin{cases}v_r,&0\le r<d,\\
       \sigma(v_{2d-1-r}),&d\le r<2d.
       \end{cases}
\]
These are precisely the letters of $(v\theta(v))^\omega$.
\end{proof}

\begin{proposition}\label{prop:sharp}
The sufficient formula $p+q-\gcd(p,q)$ cannot be replaced uniformly
by $p+q-\gcd(p,q)-1$.
\end{proposition}
\begin{proof}
Take distinct letters $a,b$, let $\theta$ be reversal, and set
$w=abb$, $p=2$, $q=3$. The word $abb$ is a prefix of
$(ab\,ba)^\omega$ and of $(abb\,bba)^\omega$, so both alternating
periods occur. Its length is $3=2+3-1-1$, but alternating period $1$
would force the constant word $aaa\cdots$. It does not occur.
\end{proof}
This example establishes only the stated uniform obstruction. It does
not assert a sharp minimum length for each pair $(p,q)$, alphabet,
or involution. Theorem~\ref{thm:main} also covers $p=q$ and divisibility
of one period by the other.

\section{Related results and verification}
The ordinary Fine--Wilf threshold and reflection/doubled-period
mechanisms are prior work \cite{FW,OWR,Bischoff}. Theorems 25 and 26
of Czeizler--Kari--Seki \cite{CKS} concern a $\theta$-palindromic
base word and common $\theta$-primitive roots. Their Theorem 28 doubles
a base block to $u\theta(u)$; Corollary 32 treats more general block
powers at the longer bound $2p+q-d$. Kari--Seki \cite{KS} improve
that general sufficient bound to $2p+q-d-\lfloor d/2\rfloor$ for
$p>q\ge2d$. These hypotheses and conclusions differ from the
alternating-prefix assertion proved here.

Simpson's published Theorem 3.4 \cite{Simpson} concerns palindromic
periodicities with half-periods $h_1,h_2$ and reflection offsets
$r_1,r_2$, and requires length at least
$2h_1+2h_2-D$, where
$D=\gcd(2(r_2-r_1),2h_1,2h_2)$; its conclusion is an ordinary
period $D$. At the common alternating phase, $h_1=p,h_2=q$ and
$D=2d$, so that condition is $2p+2q-2d$, twice the length threshold
in Theorem~\ref{thm:main}. Thus this close prior result does not
directly give the short bound on $w$.

A bounded priority audit inspected these primary results, related
proofs, bibliographies, equivalent formulations, and 53 recorded
literature queries. It found no earlier exact theorem in that scope.
This finding does not certify first discovery or universal novelty.
The accompanying audit records its specific access and version gaps,
including the unavailable Kari--Seki publisher full text and an
unrecovered earlier conference version on palindromes and local
periodicity. The cited Kari--Seki comparison uses its dated author
manuscript; the Czeizler--Kari--Seki comparison uses the actual
2010 journal-layout version, separately from the MFCS 2008 version.

The verification supplement preserves the submitted proof and
programs, three independent mathematical audit families, the bounded
priority record, complete finite-control outputs, and checksums.
The families use signed constraints and a reflection lift, direct
word-overlap checks, and definition/counterexample tests with
canonical finite models. They include fixed letters, swapped letters,
short truncations, divisibility cases, and below-threshold failures.
These computations supplement the universal proof above; they do not
replace it or certify priority. Public verification explicitly omits
private source downloads and execution histories.

\paragraph{AI use and review status.}
AI tools were used extensively in solving, proof checking, literature
research, drafting, computation, and independent adversarial agent
reviews. This is an unrefereed research note. Those automated reviews
are not independent external human peer review; no such human review
is claimed.

\begin{thebibliography}{9}
\bibitem{FW}
N. J. Fine and H. S. Wilf,
Uniqueness theorems for periodic functions,
\emph{Proc. Amer. Math. Soc.} \textbf{16}(1) (1965), 109--114.
\href{https://doi.org/10.1090/S0002-9939-1965-0174934-9}{doi:10.1090/S0002-9939-1965-0174934-9}.
\bibitem{OWR}
D. Nowotka (joint work with B. Bischoff), Word periods under involution,
in \emph{Mini-Workshop: Combinatorics on Words},
\emph{Oberwolfach Reports} \textbf{7}(3), Report 37/2010,
2195--2244 (contribution 2219--2222; published 2011).
\href{https://doi.org/10.4171/OWR/2010/37}{doi:10.4171/OWR/2010/37}.
\bibitem{Bischoff}
B. Bischoff, \emph{Wortbegrenzung unter Involution},
Diplomarbeit 3095, Universit\"at Stuttgart, 18 November 2010, Chapter 3.
\href{https://www2.informatik.uni-stuttgart.de/bibliothek/ftp/medoc_restrict.ustuttgart_fi/DIP-3095/DIP-3095.pdf}{Institutional full text}.
\bibitem{CKS}
E. Czeizler, L. Kari and S. Seki,
On a special class of primitive words,
\emph{Theoret. Comput. Sci.} \textbf{411}(3) (2010), 617--630.
\href{https://doi.org/10.1016/j.tcs.2009.09.037}{doi:10.1016/j.tcs.2009.09.037}.
\bibitem{KS}
L. Kari and S. Seki,
An improved bound for an extension of Fine and Wilf's theorem
and its optimality,
\emph{Fund. Inform.} \textbf{101}(3) (2010), 215--236.
\href{https://doi.org/10.3233/FI-2010-285}{doi:10.3233/FI-2010-285}.
\bibitem{Simpson}
J. Simpson, Palindromic periodicities,
\emph{Australas. J. Combin.} \textbf{92}(3) (2025), 450--462.
\href{https://ajc.maths.uq.edu.au/pdf/92/ajc_v92_p450.pdf}{Published full text}.
\end{thebibliography}
\end{document}
